\newpage
\setcounter{dang}{0}
% \setcounter{section}{12}
\section{Các số đặc trưng đo xu thế trung tâm}
\subsection{KIẾN THỨC TRỌNG TÂM}
\subsubsection{Số trung bình cộng (số trung bình)}
\begin{boxdn}
	Số trung bình cộng của mẫu số liệu $x_1,x_2,\ldots,x_n$, kí hiệu là $\overline{x}$, được tính bằng công thức \fbox{$\overline{x}=\dfrac{x_1+x_2+\cdots +x_n}{n}$}.
\end{boxdn}

	% \begin{enumerate}
	% 	\item 
		Số trung bình cộng của mẫu số liệu thống kê trong bảng phân bố tần số là
		\begin{center}
			\begin{tabular}{|>{\centering\arraybackslash}p{2cm}|>{\centering\arraybackslash}p{1.2cm}|>{\centering\arraybackslash}p{1.2cm}|>{\centering\arraybackslash}p{1.2cm}|>{\centering\arraybackslash}p{1.2cm}|}
				\hline
				Giá trị	& $x_1$ & $x_2$ & $\ldots$ & $x_k$ \\
				\hline
				Tần số	& $n_1$ & $n_2$ & $\ldots$ & $n_k$ \\
				\hline
			\end{tabular}
		\end{center}
		được tính theo công thức
		\fbox{$\overline{x}=\dfrac{n_1x_1 + n_2x_2 + \ldots + n_kx_k}{n},$}
		trong đó $n_i$ tương ứng là tần số của giá trị $x_i$ ($i=1;2;\ldots; k$) và $n=n_1+n_2+\cdots +n_k$.
		% \item Số trung bình cộng của mẫu số liệu thống kê trong bảng phân bố tần số tương đối
		% \begin{center}
		% 	\begin{tabular}{|>{\centering\arraybackslash}p{3.5cm}|>{\centering\arraybackslash}p{1.2cm}|>{\centering\arraybackslash}p{1.2cm}|>{\centering\arraybackslash}p{1.2cm}|>{\centering\arraybackslash}p{1.2cm}|}
		% 		\hline
		% 		Giá trị	& $x_1$ & $x_2$ & $\ldots$ & $x_k$ \\
		% 		\hline
		% 		Tần số tương đối	& $f_1$ & $f_2$ & $\ldots$ & $f_k$ \\
		% 		\hline
		% 	\end{tabular}
		% \end{center}
		% được tính theo công thức
		% \fbox{$\overline{x}=f_1x_1+f_2x_2+\cdots +f_kx_k,$}
		% trong đó $f_1=\dfrac{n_1}{n}$, $f_2=\dfrac{n_2}{n}$, $\ldots$, $f_k=\dfrac{n_k}{n}$ với $n= n_1+n_2+\cdots +n_k$.
	% \end{enumerate}

% \textbf{Ý nghĩa.}
Số trung bình là giá trị trung bình cộng của các số trong mẫu số liệu, nó cho biết vị trí trung tâm của mẫu số liệu và có thể dùng để đại diện cho mẫu số liệu.

\subsubsection{Trung vị}
\begin{boxdn}
	Sắp thứ tự mẫu số liệu gồm $n$ số liệu thành một dãy không giảm (hoặc không tăng).
	\begin{itemize}
		\item Nếu $n$ là lẻ thì số liệu đứng ở vị trí thứ $\dfrac{n+1}{2}$ gọi là trung vị.
		\item Nếu $n$ là chẵn thì số trung bình cộng của hai số liệu đứng ở vị trí thứ $\dfrac{n}{2}$ và $\dfrac{n}{2}+1$ gọi là trung vị.
	\end{itemize}
	Trung vị kí hiệu là $M_e$.
\end{boxdn}
% \begin{note}
% 	Trung vị không nhất thiết là một số trong mẫu số liệu. Khi các số liệu trong mẫu không có sự chênh lệnh lớn thì số trung bình cộng và trung vị xấp xỉ nhau.
% \end{note}
\textbf{Ý nghĩa.} Trung vị là giá trị chia đôi mẫu số liệu, nghĩa là trong mẫu số liệu được sắp xếp theo thứ tự không giảm thì giá trị trung vị ở vị trí chính giữa. Trung vị không bị ảnh hưởng bởi giá trị bất thường trong khi số trung bình bị ảnh hưởng bởi giá trị bất thường.

\subsubsection{Tứ phân vị}
\begin{boxdn}
	Sắp thứ tự mẫu số liệu gồm $n$ số liệu thành một dãy không giảm.\\
	Tứ phân vị của mẫu số liệu trên là bộ ba giá trị tứ phân vị thứ nhất, tứ phân vị thứ hai và tứ phân vị thứ ba, ba giá trị này chia mẫu số liệu thành bốn phần có số phần tử bằng nhau.
	\begin{itemize}
		\item Tứ phân vị $Q_2$ bằng trung vị.
		\item Nếu $n$ là chẵn thì tứ phân vị thứ nhất $Q_1$ bằng trung vị của nửa dãy phía dưới và tứ phân vị thứ ba $Q_3$ bằng trung vị của nửa dãy phía trên.
		\item Nếu $n$ là lẻ thì tứ phân vị thứ nhất $Q_1$ bằng trung vị của nửa dãy phía dưới (không bao gồm $Q_2$) và tứ phân vị thứ ba $Q_3$ bằng trung vị của nửa dãy phía trên (không bao gồm $Q_2$).
	\end{itemize}
	\begin{center}
		% \begin{tikzpicture}[scale=0.9, font=\footnotesize, line join=round, line cap=round, >=stealth]
		% 	\path (0,0) coordinate(A) ++(0:2) coordinate(Q_1) ++(0:3.5) coordinate(Q_2) ++(0:2.5) coordinate(Q_3) ++(0:1.5) coordinate(B);
		% 	\draw (A)node[text width=1.5cm,align=center,below=5pt] {Giá trị\\nhỏ nhất} -- (Q_1) node[below=5pt]{$Q_1$} -- (Q_2) node[below=5pt]{$Q_2$} -- (Q_3) node[below=5pt]{$Q_3$} -- (B)node[below=5pt, text width=1.5 cm, align=center]{Giá trị\\lớn nhất};
		% 	\foreach \i in {A,Q_1,Q_2,Q_3,B} \draw (\i)++(-90:0.07) -- ++(90:0.14);
		% 	\draw[decoration={brace}, decorate] ($(A)+(0,0.2)$)--($(Q_2)+(0,0.2)$) node[sloped,midway,above]{Nửa dãy phía dưới};
		% 	\draw[decoration={brace}, decorate] ($(Q_2)+(0,0.2)$)--($(B)+(0,0.2)$) node[sloped,midway,above]{Nửa dãy phía trên};	
		% \end{tikzpicture}
		\begin{tikzpicture}[scale=0.9, font=\footnotesize, line join=round, line cap=round, >=stealth]
				\path (0,0) coordinate(A) ++(0:2) coordinate(Q_1) ++(0:3.5) coordinate(Q_2) ++(0:2.5) coordinate(Q_3) ++(0:1.5) coordinate(B);
				\draw (A)node[text width=1.5cm,align=center,below=5pt] {Giá trị\\nhỏ nhất} -- (Q_1) node[below=5pt]{$Q_1$} -- (Q_2) node[below=5pt]{$Q_2$} -- (Q_3) node[below=5pt]{$Q_3$} -- (B)node[below=5pt, text width=1.5 cm, align=center]{Giá trị\\lớn nhất};
				\foreach \i in {A,Q_1,Q_2,Q_3,B} \draw (\i)++(-90:0.07) -- ++(90:0.14);
				\draw[decoration={brace}, decorate] ($(A)+(0,0.2)$)--($(Q_1)+(0,0.2)$) node[sloped,midway,above]{$25\%$};
				\draw[decoration={brace,mirror}, decorate,yscale=-1] ($(Q_1)+(0,0.2)$)--($(Q_2)+(0,0.2)$) node[sloped,midway,below]{$25\%$};
				\draw[decoration={brace,mirror}, decorate,yscale=-1] ($(Q_2)+(0,0.2)$)--($(Q_3)+(0,0.2)$) node[sloped,midway,below]{$25\%$};	
				\draw[decoration={brace}, decorate] ($(Q_3)+(0,0.2)$)--($(B)+(0,0.2)$) node[sloped,midway,above]{$25\%$};
			\end{tikzpicture}
	\end{center}
\end{boxdn}
% \begin{note}\,
% 	\begin{itemize}
% 		\item $Q_1$ còn được gọi là tứ phân vị dưới, $Q_3$ còn được gọi là tứ phân vị trên.
% 		\item \textbf{Ý nghĩa.} Các điểm $Q_1$, $Q_2$, $Q_3$ chia dãy dữ liệu đã sắp xếp theo thứ tự từ nhỏ đến lớn thành $4$ phần, mỗi phần đều chứa $25\%$ giá trị.
% 		\begin{center}
% 			\begin{tikzpicture}[scale=0.9, font=\footnotesize, line join=round, line cap=round, >=stealth]
% 				\path (0,0) coordinate(A) ++(0:2) coordinate(Q_1) ++(0:3.5) coordinate(Q_2) ++(0:2.5) coordinate(Q_3) ++(0:1.5) coordinate(B);
% 				\draw (A)node[text width=1.5cm,align=center,below=5pt] {Giá trị\\nhỏ nhất} -- (Q_1) node[below=5pt]{$Q_1$} -- (Q_2) node[below=5pt]{$Q_2$} -- (Q_3) node[below=5pt]{$Q_3$} -- (B)node[below=5pt, text width=1.5 cm, align=center]{Giá trị\\lớn nhất};
% 				\foreach \i in {A,Q_1,Q_2,Q_3,B} \draw (\i)++(-90:0.07) -- ++(90:0.14);
% 				\draw[decoration={brace}, decorate] ($(A)+(0,0.2)$)--($(Q_1)+(0,0.2)$) node[sloped,midway,above]{$25\%$};
% 				\draw[decoration={brace,mirror}, decorate,yscale=-1] ($(Q_1)+(0,0.2)$)--($(Q_2)+(0,0.2)$) node[sloped,midway,below]{$25\%$};
% 				\draw[decoration={brace,mirror}, decorate,yscale=-1] ($(Q_2)+(0,0.2)$)--($(Q_3)+(0,0.2)$) node[sloped,midway,below]{$25\%$};	
% 				\draw[decoration={brace}, decorate] ($(Q_3)+(0,0.2)$)--($(B)+(0,0.2)$) node[sloped,midway,above]{$25\%$};
% 			\end{tikzpicture}
% 		\end{center}
% 	\end{itemize}
% \end{note}
\textbf{Ý nghĩa.} Các điểm $Q_1$, $Q_2$, $Q_3$ chia dãy dữ liệu đã sắp xếp theo thứ tự từ nhỏ đến lớn thành $4$ phần, mỗi phần đều chứa $25\%$ số giá trị.
\subsubsection{Mốt}
\begin{boxdn}
	Mốt của mẫu số liệu là giá trị có tần số lớn nhất trong bảng phân bố tần số và kí hiệu là $M_\circ$.
\end{boxdn}
\begin{note}
	Một mẫu số liệu có thể có một hoặc nhiều mốt
\end{note}

%%%%%%
\subsection{PHÂN LOẠI VÀ PHƯƠNG PHÁP GIẢI TOÁN}

\begin{dang}{Số trung bình}
	\underline{Phương pháp:} $\bar{x} = \dfrac{\text{Tổng các giá trị}}{\text{Số các giá trị}}$
\end{dang}
\begin{vd}%[0D6N3-2]
	Trong một cuộc thi tìm hiểu lịch sử địa phương (thang điểm $10$), một lớp học tham gia cuộc thi và đạt được số điểm như sau
	\begin{center}
		\begin{tabular}{|c|c|c|c|c|}
			\hline
			Số học sinh & $5$ & $12$ & $10$ & $3$\\
			\hline
			Số điểm & $5$ & $6$ & $7$ & $9$\\
			\hline
		\end{tabular}
	\end{center}
	Hỏi trung bình mỗi học sinh trong lớp đạt bao nhiêu điểm trong cuộc thi?
	\loigiai{
		\begin{itemize}
			\item Số học sinh của lớp là $5+12+10+3=30$ học sinh.
			\item Điểm trung bình mà mỗi bạn đạt được trong cuộc thi là \\
			\centerline{$\dfrac{5\cdot5 + 12\cdot6 + 10\cdot7 + 3\cdot9}{30}\approx 6{,}47$ (điểm).}
		\end{itemize}
	}
\end{vd}

\begin{vd}%[0D6N3-2]
	Khi nghiên cứu tuổi thọ của một loại bóng đèn, người ta đã chọn tùy ý $10$ bóng đèn trong một lô hàng và bật sáng liên tục cho đến khi nó tự tắt. Tuổi thọ của các bóng đèn (tính theo giờ) của các bóng đèn được ghi lại trong bảng sau
	\begin{center}
		\begin{tabular}{|c|c|c|c|c|}
			\hline
			Số bóng đèn & $2$ & $3$ & $4$ & $1$\\
			\hline
			Tuổi thọ (giờ) & $1\,150$ & $1\,160$ & $1\,170$ & $1\,180$\\
			\hline
		\end{tabular}
	\end{center}
	Hỏi tuổi thọ trung bình của các bóng đèn trong lô hàng là bao nhiêu?
	\loigiai{
		Tuổi thọ trung bình của các bóng đèn trong lô hàng là \\[.5em]
		\centerline{$\dfrac{2\cdot1\,150 + 3\cdot1\,160 + 4\cdot1\,170 + 1\cdot1\,180}{10}=1\,164$ (giờ).}
	}
\end{vd}


\begin{vd}%[0D6N3-2]
	Trong đợt kiểm tra quân sự thường niên tại một đơn vị, ở bộ môn bắn súng AK mỗi người phải bắn $5$ phát súng vào bia. Thang điểm bắn là $0$, $4$, $5$, $6$, $7$, $8$, $9$, $10$. Ở $4$ lần bắn trước đó, anh Nam đã đạt được số điểm như sau
	\begin{center}
		\begin{tabular}{|c|c|c|c|c|}
			\hline
			Lần bắn & $1$ & $2$ & $3$ & $4$ \\
			\hline
			Số điểm & $8$ & $7$ & $0$ & $9$  \\
			\hline
	\end{tabular}
	\end{center}
	Biết rằng để vượt qua bài kiểm tra, mỗi người phải đạt điểm trung bình trong các lần bắn từ $6{,}5$ điểm trở lên. Tính số điểm ít nhất mà anh Nam cần đạt được trong lần bắn thứ $5$ để vượt qua bài kiểm tra.
	\loigiai{
		\begin{itemize}
			\item Gọi điểm số trong lần bắn thứ 5 của anh Nam là $n$ điểm, $n \in \{0,4,5,6,7,8,9,10\}$.\\
			Khi đó điểm số trung bình mà anh Nam đạt được trong $5$ lần bắn là\\[.5em]
			\centerline{$\dfrac{8+7+0+9+n}{5}$}.
			\item Theo yêu cầu đề bài, ta có $\dfrac{8+7+0+9+n}{5} \ge 6{,}5 \Leftrightarrow n \ge 8{,}5$.\\
			Vì $n \in \{0,4,5,6,7,8,9,10\}$ nên số điểm ít nhất mà anh Nam cần đạt được trong lần bắn thứ $5$ để vượt qua bài điểm tra là $9$ điểm.
		\end{itemize}
	}
\end{vd}
\begin{dang}{Số trung vị}
	Sắp thứ tự mẫu số liệu gồm $n$ số liệu thành một dãy không giảm (hoặc không tăng).\\
	Nếu $n$ là lẻ thì $M_e=\text{Giá trị đứng chính giữa}$.\\
	Nếu $n$ là chẵn thì $M_e=\text{Trung bình cộng của hai giá trị đứng chính giữa}$.
\end{dang}
\begin{vd}%[0D6H3-3]
	Một nhóm gồm 7 học sinh tham gia một cuộc thi và đạt được số điểm như sau: $89$, $69$, $65$, $0$, $80$, $0$, $90$.
	Hãy tìm trung vị của mẫu số liệu trên.
	\loigiai{
		\begin{itemize}
			\item Sắp xếp dãy số liệu theo thứ tự không giảm như sau $0$, $0$, $65$, $69$, $80$, $89$, $90$.
			\item Dãy trên có $7$ giá trị, giá trị chính giữa là $69$. Vậy trung vị của dãy số liệu trên bằng $69$.
		\end{itemize}
	}
\end{vd}

\begin{vd}%[0D6H3-3]
	Số áo bán được trong một cửa hàng trong một quý được ghi lại trong bảng sau
	\begin{center}
		\begin{tabular}{|c|c|c|c|c|c|c|c|}
			\hline
			Cỡ số & $36$ & $37$ & $38$ & $39$ & $40$ & $41$ & $42$\\
			\hline
			Số áo bán được & $13$ & $45$ & $126$ & $110$ & $126$ & $40$ & $5$\\
			\hline
		\end{tabular}
	\end{center}
	Hãy tìm trung vị của mẫu số liệu trên.
	\loigiai{
		\begin{itemize}
			\item Số giá trị của dãy số liệu trên là $13+45+126+110+126+40+5=465$.
			\item Sắp xếp dãy số liệu theo thứ tự không giảm ta được\\
			$36$, $\ldots$, $36$, $37$, $\ldots$, $37$, $38$, $\ldots$, $38$, $39$, $\ldots$, $39$, $40$, $\ldots$, $40$, $41$, $\ldots$, $41$, $42$, $\ldots$, $42$.
			\item Dãy trên có $465$ giá trị, giá trị chính giữa là giá trị thứ $233$, giá trị này bằng $39$. \\
			Vậy trung vị của dãy số liệu trên bằng $39$.
		\end{itemize}
	}
\end{vd}

\begin{vd}%[0D6H3-3]
	Số đôi giày bán được trong một cửa hàng trong một tháng được ghi lại trong bảng sau
	\begin{center}
		\begin{tabular}{|c|c|c|c|c|c|c|}
			\hline
			Cỡ giày & $38$ & $39$ & $40$ & $41$ & $42$ & $43$\\
			\hline
			Số đôi bán được & $18$ & $52$ & $95$ & $110$ & $70$ & $25$\\
			\hline
		\end{tabular}
	\end{center}
	Hãy tìm trung vị của mẫu số liệu trên.
	\loigiai{
		\begin{itemize}
			\item Số giá trị của mẫu số liệu là $18+52+95+110+70+25=370$.
			\item Sắp xếp dãy số liệu theo thứ tự không giảm ta được\\
			$38$, $\ldots$, $38$, $39$, $\ldots$, $39$, $40$, $\ldots$, $40$, $41$, $\ldots$, $41$, $42$, $\ldots$, $42$, $43$, $\ldots$, $43$.
			\item Vì dãy số liệu có $370$ giá trị (là số chẵn) nên trung vị là trung bình cộng của hai giá trị thứ $185$ và $186$.\\
			Hai giá trị thứ $185$ và $186$ đều thuộc cỡ giày $41$.
			\item Vậy trung vị của mẫu số liệu trên bằng $41$.
		\end{itemize}
	}
\end{vd}


\begin{dang}{Tứ phân vị}
	\underline{Phương pháp:} Tứ phân vị $Q_2$ bằng trung vị, tứ phân vị thứ nhất $Q_1$ bằng trung vị của nửa dãy phía dưới và tứ phân vị thứ ba $Q_3$ bằng trung vị của nửa dãy phía trên\\
	\underline{Chú ý:} Nếu $n$ là lẻ mỗi nửa dãy trên và dưới không bao gồm $Q_2$.
	\begin{center}
		\begin{tikzpicture}[scale=0.9, font=\footnotesize, line join=round, line cap=round, >=stealth]
			\path (0,0) coordinate(A) ++(0:2) coordinate(Q_1) ++(0:3.5) coordinate(Q_2) ++(0:2.5) coordinate(Q_3) ++(0:1.5) coordinate(B);
			\draw (A)node[text width=1.5cm,align=center,below=5pt] {Giá trị\\nhỏ nhất} -- (Q_1) node[below=5pt]{$Q_1$} -- (Q_2) node[below=5pt]{$Q_2$} -- (Q_3) node[below=5pt]{$Q_3$} -- (B)node[below=5pt, text width=1.5 cm, align=center]{Giá trị\\lớn nhất};
			\foreach \i in {A,Q_1,Q_2,Q_3,B} \draw (\i)++(-90:0.07) -- ++(90:0.14);
			\draw[decoration={brace}, decorate] ($(A)+(0,0.2)$)--($(Q_2)+(0,0.2)$) node[sloped,midway,above]{Nửa dãy phía dưới};
			\draw[decoration={brace}, decorate] ($(Q_2)+(0,0.2)$)--($(B)+(0,0.2)$) node[sloped,midway,above]{Nửa dãy phía trên};	
		\end{tikzpicture}
	\end{center}
\end{dang}
\begin{vd}%[0D6H3-4]
	Số tấn hàng bán ra được trong $6$ tháng đầu năm của một công ty được cho như sau $4$, $7$, $9$, $11$, $12$, $20$. Tìm tứ phân vị dưới của mẫu số liệu trên.
	\loigiai{
		\begin{itemize}
			\item Sắp xếp mẫu số liệu trên theo thứ tự không giảm $4$, $7$, $9$, $11$, $12$, $20$.
			\item Mẫu số liệu trên có $6$ giá trị, nên trung vị là trung bình cộng của hai số hạng chính giữa $Q_2=\dfrac{9+11}{2}=10$.
			\item Như vậy, nửa số liệu bên trái $Q_2$ là $4$, $7$, $9$. Mẫu số liệu này có $3$ giá trị, nên trung vị là số hạng chính giữa, vậy $Q_1=7$.
			\item Vậy tứ phân vị dưới của mẫu số liệu đã cho là $Q_1=7$.
		\end{itemize}
	}
\end{vd}


\begin{vd}%[0D6H3-4]
	Số giờ tự học trong một tuần của một nhóm học sinh được cho như sau:
	$4$, $6$, $9$, $17$, $20$, $12$, $14$.
	Tìm tứ phân vị trên và tứ phân vị dưới của mẫu số liệu đã cho.
	\loigiai{
		\begin{itemize}
			\item Sắp xếp mẫu số liệu theo thứ tự không giảm là
			$4$, $6$, $9$, $12$, $14$, $17$, $20$.
			\item Mẫu số liệu có $7$ giá trị nên trung vị là số hạng chính giữa là $Q_2 = 12$.
			\item Nửa số liệu bên trái của $Q_2$ là $4$, $6$, $9$.
			Mẫu số liệu này có $3$ giá trị nên trung vị là số hạng chính giữa, do đó $Q_1 = 6$.
			\item Nửa số liệu bên phải của $Q_2$ là
			$14$, $17$, $20$.\\
			Mẫu số liệu này có $3$ giá trị nên trung vị là số hạng chính giữa, do đó $Q_3 = 17$.
			\item Vậy tứ phân vị dưới và tứ phân vị trên của mẫu số liệu đã cho lần lượt là
			$$
			Q_1 = 6,\quad Q_3 = 17.
			$$
		\end{itemize}
	}
\end{vd}

\begin{vd}
Số đôi giày bán được trong một cửa hàng trong một tháng được ghi lại trong bảng sau
	\begin{center}
		\begin{tabular}{|c|c|c|c|c|c|c|}
			\hline
			Cỡ giày & $38$ & $39$ & $40$ & $41$ & $42$ & $43$\\
			\hline
			Số đôi bán được & $18$ & $52$ & $95$ & $110$ & $70$ & $25$\\
			\hline
		\end{tabular}
	\end{center}
	Hãy tìm tứ phân vị thứ nhất và thứ ba của mẫu số liệu trên.
\end{vd}
%%%%%%%%%%%%%
%%%%%%%%%%%%%
%%%%%%%%%%%%%
\begin{dang}{Mốt}
	Mốt của mẫu số liệu là giá trị có tần số lớn nhất trong bảng phân bố tần số và kí hiệu là $M_\circ$.
\end{dang}
\begin{vd}%[0D6N3-5]
	Giá thành của một sản phẩm (tính theo đơn vị nghìn đồng) của $20$ cơ sở sản xuất được cho bởi bảng sau
	\begin{center}
		\begin{tabular}{|c|c|c|c|c|c|c|c|c|c|}
			\hline
			15 & 25 & 25 & 30 & 20 & 25 & 35 & 30 & 25 & 30\\
			\hline
			25 & 20 & 35 & 30 & 15 & 25 & 25 & 20 & 25 & 25\\
			\hline
		\end{tabular}
	\end{center}
	Tìm mốt của mẫu số liệu trên.
	\loigiai{
		\begin{itemize}
			\item Lập bảng tần số (số lần xuất hiện của giá trị) cho mẫu số liệu đã cho, ta được
			\begin{center}
				\begin{tabular}{|c|c|c|c|c|c|}
					\hline
					Giá trị & 15 & 20 & 25 & 30 & 35\\
					\hline
					Tần số & 2 & 3 & 9 & 4 & 2\\
					\hline
				\end{tabular}
			\end{center}
			\item Dựa vào bảng tần số, ta thấy giá trị $25$ xuất hiện nhiều nhất ($9$ lần) nên mốt của dấu hiệu là $25$.
		\end{itemize}
	}
\end{vd}
%%%%%%%%%%%%%
\begin{vd}%[0D6N3-5]
	Số cân nặng của $20$ học sinh được ghi lại như sau
	\begin{center}
		\begin{tabular}{|c|c|c|c|c|c|c|c|c|c|}
			\hline
			28 & 35 & 29 & 37 & 30 & 35 & 37 & 30 & 35 & 29\\
			\hline
			30 & 37 & 35 & 35 & 42 & 28 & 35 & 29 & 37 & 20\\
			\hline
		\end{tabular}
	\end{center}
	Tìm mốt của mẫu số liệu trên.
	\loigiai{
		\begin{itemize}
			\item Lập bảng tần số (số lần xuất hiện của giá trị) cho mẫu số liệu đã cho, ta được
			\begin{center}
				\begin{tabular}{|c|c|c|c|c|c|c|c|}
					\hline
					Giá trị & 20 & 28 & 29 & 30 & 35 & 37 & 42\\
					\hline
					Tần số & 1 & 2 & 3 & 3 & 6 & 4 & 1\\
					\hline
				\end{tabular}
			\end{center}
			\item Dựa vào bảng tần số, ta thấy giá trị $35$ xuất hiện nhiều nhất ($6$ lần) nên mốt của dấu hiệu là $35$.
		\end{itemize}
	}
\end{vd}
%%%%%%%%%%%%%
% \begin{vd}%[0D6N3-5]
% 	Số điểm kiểm tra Toán của $20$ học sinh trong một lớp được ghi lại như sau
% 	\begin{center}
% 		\begin{tabular}{|c|c|c|c|c|c|c|c|c|c|}
% 			\hline
% 			6 & 7 & 8 & 7 & 9 & 8 & 7 & 6 & 8 & 7\\
% 			\hline
% 			8 & 9 & 7 & 6 & 8 & 7 & 9 & 8 & 7 & 6\\
% 			\hline
% 		\end{tabular}
% 	\end{center}
% 	Tìm mốt của mẫu số liệu trên.
% 	\loigiai{
% 		\begin{itemize}
% 			\item Lập bảng tần số (số lần xuất hiện của mỗi giá trị) cho mẫu số liệu đã cho:
% 			\begin{center}
% 				\begin{tabular}{|c|c|c|c|c|}
% 					\hline
% 					Giá trị & 6 & 7 & 8 & 9\\
% 					\hline
% 					Tần số & 4 & 7 & 6 & 3\\
% 					\hline
% 				\end{tabular}
% 			\end{center}
% 			\item Dựa vào bảng tần số, ta thấy giá trị $7$ xuất hiện nhiều nhất ($7$ lần), nên mốt của mẫu số liệu là $7$.
% 		\end{itemize}
% 	}
% \end{vd}

\subsection{BÀI TẬP RÈN LUYỆN}
\ind{PHẦN I.} \inden{Câu trắc nghiệm nhiều phương án lựa chọn. Mỗi câu hỏi học sinh chỉ chọn một phương án.}
\setcounter{ex}{0}
\Opensolutionfile{ans}[ans/ans-TN-C5-B3-D1]
%%%%===== Câu 1
\begin{ex}%[0D6N3-2]
	Tuổi đời của $16$ công nhân trong xưởng sản xuất được thống kê trong bảng sau
	\begin{center}
		\begin{tabular}{|c|c|c|c|c|c|c|c|}
			\hline Tuổi & $25$ & $26$ & $27$ & $29$ & $30$ & $33$ & Cộng \\
			\hline Số người & $2$ & $3$ & $4$ & $3$ & $3$ & $1$ & $16$ \\
			\hline
		\end{tabular}
	\end{center}
	Tìm số trung bình $\overline{x}$ của mẫu số liệu trên
	\choice
	{$28$}
	{$27{,}75$}
	{\True $27{,}875$}
	{$27$}
	\loigiai{		
		Mẫu số liệu có số trung bình là 
		\allowdisplaybreaks
		\begin{eqnarray*}
			\overline x =\dfrac{25\cdot 2+26\cdot 3+27\cdot 4+29\cdot 3+30\cdot 3+33\cdot 1}{16}= 27{,}875.
		\end{eqnarray*}
	}
\end{ex}

%%%%===== Câu 2
\begin{ex}%[0D6N3-3]
	Cho mẫu thống kê $\left\{28,\,16,\,13,\,18,\,12,\,28,\,13,\,19\right\}$ . Trung vị của mẫu số liệu trên là bao nhiêu?
	\choice
	{$14$}
	{\True $17$}
	{$18$}
	{$20$}
	\loigiai{		
		Mẫu thống kê trên có $ 8$ số liệu được sắp xếp theo thứ tự không giảm là $12,\,13,\,13,\,16,\,18,\,19,\,28,\,28$, nên trung vị của mẫu số liệu trên là $M_e=\dfrac{16+18}{2}=17$.}
\end{ex}

%%%%===== Câu 3
\begin{ex}%[0D6N3-5]
	Cho mẫu số liêu thống kê $\{6, 5, 5, 2, 9, 10, 8\}$. Mốt của mẫu số liêu trên bằng bao nhiêu?
	\choice
	{\True $5$}
	{$10$}
	{$2$}
	{$6$}
	\loigiai{
		Mốt của mẫu số liêu trên là $M_o=5$.
	}
\end{ex}

%%%%===== Câu 4
\begin{ex}%[0D6N3-3]
	Cho bảng số liệu thống kê chiều cao của một nhóm học sinh như sau
\begin{center}
	\begin{tabular}{|l|l|l|l|l|l|l|l|}
		\hline
		$150$ & $153$ & $153$ & $154$ & $154$ & $155$ & $160$ & $160$ \\
		\hline
		$162$ & $162$ & $163$ & $163$ & $163$ & $165$ & $165$ & $167$ \\
		\hline
	\end{tabular}
\end{center}
	Số trung vị của bảng số liệu nói trên là
	\choice
	{\True $161$}
	{$153$}
	{$163$}
	{$156$}
	\loigiai{		
		Ta có trong bảng số liệu thống kê có tất cả $ 16$ giá trị. Do đó số trung vị bằng trung bình cộng của hai số đứng thứ $ 8$ và $ 9$ trong bảng số liệu thống kê.\\
		Ta có $M_e=\dfrac{160+162}{2}=161.$}
\end{ex}

%%%%===== Câu 5
\begin{ex}%[0D6N3-5]
	Cho bảng phân bố tần số tiền thưởng (triệu đồng) cho cán bộ và nhân viên trong một công ty\\
	\centerline{\begin{tabular}{|c|c|c|c|c|c|c|}
			\hline
			Tiền thưởng & $2$ & $3$ & $4$ & $5$ & $6$ & Cộng\\
			\hline
			Tần số & $5$ & $15$ & $10$ & $6$ & $7$ & $43$\\
			\hline
	\end{tabular}}\\
	Mốt của bảng phân bố tần số đã cho là
	\choice
	{\True $3$ triệu đồng}
	{$2$ triệu đồng}
	{$6$ triệu đồng}
	{$5$ triệu đồng}
	\loigiai{
		Mốt của bảng phân bố tần số là giá trị (xi) có tần số (ni) lớn nhất và được kí hiệu là $M_o$.}
\end{ex}


%%%%===== Câu 6
\begin{ex}%[0D6N3-5]
	Mốt của một bảng phân bố tần số là
	\choice
	{tần số lớn nhất trong bảng phân bố tần số}
	{\True giá trị có tần số lớn nhất trong bảng phân bố tần số}
	{giá trị có tần số nhỏ nhất trong bảng phân bố tần số}
	{tần số nhỏ nhất trong bảng phân bố tần số}
	\loigiai{
		Mốt của một bảng phân bố tần số là giá trị có tần số lớn nhất.}
\end{ex}

%%%%===== Câu 7
\begin{ex}%[0D6H3-3]
	Cho bảng số liệu ghi lại điểm của $ 40$ học sinh trong bài kiểm tra $ 1$ tiết môn Toán.
	\begin{center}
		\begin{tabular}{|c|c|c|c|c|c|c|c|c|c|}
			\hline
			Điểm & $ 3$ & $ 4$ & $ 5$ & $ 6$ & $ 7$ & $ 8$ & $ 9$ & $ 10$ & Cộng\\
			\hline
			Số học sinh & $ 2$ & $ 3$ & $ 7$ & $18$ & $ 3$ & $ 2$ & $ 4$ & $ 1$ & $ 40$\\
			\hline
		\end{tabular}
	\end{center}
	Số trung vị là
	\choice
	{$5$}
	{\True $6$}
	{$6{,}5$}
	{$7$}
	\loigiai{
		Số trung vị là $M_e=\dfrac{x_{20}+x_{21}}{2}=\dfrac{6+6}{2}=6$.}
\end{ex}

%%%%===== Câu 8
\begin{ex}%[0D6H3-5]
	Một cửa hàng trà sữa vừa khai trương, thống kê lượng khách tới quán trong $7$ ngày đầu và thu được mẫu số liệu sau
	\begin{center}
		\begin{tabular}{|c|c|c|c|c|c|c|}
			\hline
			Ngày $1$ & Ngày $2$ & Ngày $3$ & Ngày $4$ & Ngày $5$ & Ngày $6$ & Ngày $7$\\
			\hline
			$575$ & $454$ & $400$ & $325$ & $351$ & $333$ & $412$\\
			\hline
		\end{tabular}
	\end{center}
	Mệnh đề nào sau đây là đúng?
	\choice
	{Số trung vị là $263$}
	{\True Số trung bình làm tròn đến hàng phần trăm là $407{,}14$}
	{Dấu hiệu điều tra ở đây là doanh thu của quán trà sữa}
	{Ngày $2$ là mốt của mẫu số liệu này}
	\loigiai{
		Ta có $\overline x=\dfrac{575+454+400+325+351+333+412}{7}=407{,}14$.}
\end{ex}

%%%%===== Câu 9
\begin{ex}%[0D6H3-2]
	Điểm trung bình thi học kỳ II môn Toán của một nhóm gồm $N$ học sinh lớp 12A6 là $8{,}1$ . Biết rằng tổng điểm môn toán của nhóm này là $72{,}9$. Tìm số học sinh của nhóm.
	\choice
	{$20$}
	{\True $9$}
	{$8$}
	{$15$}
	\loigiai{		
		Ta có giá giá trị $N=\dfrac{72{,}9}{8{,}1}=9$ (học sinh).}
\end{ex}

%%%%===== Câu 10
\begin{ex}%[0D6H3-3]
	Kết quả điểm kiểm tra môn Toán trong một kì thi của $ 200$ em học sinh được trình bày ở bảng sau
	\begin{center}
		\begin{tabular}{|c|c|c|c|c|c|c|c|}
			\hline Điềm & $5$ & $6$ & $7$ & $8$ & $9$ & $10$ & Cộng \\
			\hline Tần số & $10$ & $35$ & $38$ & $63$ & $42$ & $12$ & $200$ \\
			\hline
		\end{tabular}
	\end{center}
	Số trung vị của bản phân bố tần số nói trên là
	\choice
	{\True $8$}
	{$7$}
	{$6$}
	{$9$}
	\loigiai{
		Số trung vị của bản phân bố tần số nói trên là $8$.}
\end{ex}


%%%%===== Câu 11
\begin{ex}%[0D6H3-3]
	Điểm thi học kì của một học sinh như sau: $4;\,6;\,2;\,7;\,3;\,5;\,9;\,8;\,7;\,10;\,9$. Số trung bình và số trung vị lần lượt là
	\choice
	{$6{,}22$ và $7$}
	{$7$ và $6$}
	{\True $6{,}36$ và $7$}
	{$6$ và $6$}
	\loigiai{
		Mẫu thống kê trên có $11$ số liệu được sắp xếp theo thứ tự không giảm là $2;\,3;\,4;\,5;\,6;\,7;\,7;\,8;\,9;\,9;\,10$.\\
		Số trung bình là $\overline x=\dfrac{2+3+4+5+6+2\cdot 7+8+2\cdot 9+10}{11}=\dfrac{70}{11}\approx 6{,}36$.\\
		Trung vị của mẫu số liệu trên là $M_e=7$.
	}
\end{ex}

%%%%===== Câu 12
\begin{ex}%[0D6H3-4]
	Cho mẫu số liệu $21; 35; 17; 43; 8; 59; 72 ;119$. Tứ phân vị thứ nhất, thứ hai, thứ ba lần lượt là
	\choice
	{ \True $19; \,39; \,65{,}5$}
	{ $26; \,43; \,65{,}5$}
	{$39; \,19; \,65{,}5$}
	{$43; \,26;\, 65{,}5$}
	\loigiai{
		Sắp xếp lại mẫu số liệu theo thứ tự không giảm, ta được: $8; 17; 21; 35; 43; 59; 72; 119$.
		\begin{itemize}
			\item Vì cỡ mẫu là $n=8$, là số chẵn, nên giá trị tứ phân vị thứ hai là
			$$
			Q_{2}=\dfrac{1}{2}(35+43)=39.
			$$
			\item Tứ phân vị thứ nhất là trung vị của mẫu: $8; 17; 21; 35$. Do đó $Q_{1}=19$.
			\item Tứ phân vị thứ ba là trung vị của mẫu:  $43; 59; 72; 119$. Do đó $Q_{3}=65{,}5$.
		\end{itemize}
	}	
\end{ex}
\Closesolutionfile{ans}
\indapan{6}{ans/ans-TN-C5-B3-D1}
%%%%%%%%%%%%%%%%%%%
\ind{PHẦN II.} \inden{Câu trắc nghiệm đúng sai. Trong mỗi ý a), b), c), d) ở mỗi câu, học sinh chọn đúng hoặc sai.}
\setcounter{ex}{0}
\Opensolutionfile{ans}[ans/ans-TF-C5-B3-D1]
%%%========Câu 1
\begin{ex}%[0D6N3-5]
	Cho mẫu số liệu sau: $4$; $5$; $6$; $7$; $8$; $4$; $9$; $4$; $3$. Khi đó
	\choiceTF
	{\True Số trung bình là $\overline x=5{,}5$}
	{Mốt là $M_o=3$}
	{Trung vị là $M_e=4$}
	{\True Tứ phân vị thứ ba là $Q_3=7$}
	\loigiai{
		\begin{itemchoice} 
			\itemch Số trung bình là  $\overline x=\dfrac{55}{10}=5{,}5$. 
			\itemch Mốt $M_o=4$.
			\itemch 	Sắp xếp mẫu $ 3; 4; 4; 4; 5; 5; 6; 7; 8; 9$. \\
			Kích thước mẫu là 10 (chẵn) nên trung vị là $M_e=\dfrac{1}{2}(5+5)=5$. 
			\itemch Tứ phân vị thứ ba là $Q_3=7$.
		\end{itemchoice}
	}
\end{ex}

%%%========Câu 2
\begin{ex}%[0D6N3-3]
	Thống kê chiều cao (đơn vị cm) của nhóm $15$ bạn nam lớp $10$ cho kết quả như sau: $162$, $157$, $170$, $165$, $166$, $157$, $159$, $164$, $172$, $155$, $156$, $156$, $180$, $165$, $155$
	Khi đó
	\choiceTF
	{Chiều cao thấp nhất là $156$}
	{\True $Q_2=162$}
	{$Q_1=157$}
	{$Q_3=170$}
	\loigiai{
		Sắp xếp số liệu theo thứ tự tăng dần ta được\\
		\centerline{\begin{tabular}{|c|c|c|c|c|c|c|c|c|c|c|c|c|c|c|}
				\hline
				$155$ & $155$ & $156$ & $156$ & $157$ & $157$ & $159$ & $162$ & $164$ & $165$ & $165$ & $166$ & $170$ & $172$ & $180$\\
				\hline
		\end{tabular}}
		\begin{itemchoice} 
			\itemch Chiều cao thấp nhất là $155$.
			\itemch Vì có $15$ giá trị nên số trung vị là số ở vị trí thứ $8$ là ${Q_2}=162$.
			\itemch 	Nửa số liệu bên trái $Q_2$ là\\
			\centerline{\begin{tabular}{|c|c|c|c|c|c|c|}
					\hline
					$155$ & $155$ & $156$ & $156$ & $157$ & $157$ & $159$\\
					\hline
			\end{tabular}}\\
			Ta tìm được trung vị $Q_1=156$.
			\itemch 	Nửa số liệu bên phải $Q_2$ là\\
			\centerline{\begin{tabular}{|c|c|c|c|c|c|c|}
					\hline
					$164$ & $165$ & $165$ & $166$ & $170$ & $172$ & $180$\\
					\hline
			\end{tabular}}\\
			Ta tìm được trung vị $Q_3=166$.
		\end{itemchoice}		
	}
\end{ex}

%%%========Câu 3
\begin{ex}%[0D6H3-4]
	Thống kê số bao xi măng được bán ra tại một cửa hàng vật liệu xây dựng trong $24$ tháng cho kết quả như sau\\
	\centerline{\begin{tabular}{|c|c|c|c|c|c|c|c|}
			\hline
			$72$ & $89$ & $88$ & $73$ & $63$ & $265$ & $69$ & $65$\\
			\hline
			$94$ & $80$ & $81$ & $98$ & $66$ & $71$ & $84$ & $73$\\
			\hline
			$93$ & $59$ & $60$ & $61$ & $83$ & $72$ & $85$ & $66$\\
			\hline
	\end{tabular}}\\
	Khi đó
	\choiceTF
	{\True  Mỗi tháng cửa hàng bán trung bình $83{,}75$ bao}
	{ Số trung vị là $72$ }
	{\True  Sai khác giữa số trung bình và số trung vị là $10{,}75$}
	{Khoảng cách từ $Q_1$ đến $Q_2$ là $8$}
	\loigiai{
		Sắp xếp lại mẫu dữ liệu theo thứ tự tăng dần ta được
		\begin{center}
			\begin{tabular}{|c|c|c|c|c|c|c|c|}
				\hline
				$59$ & $60$ & $61$ & $63$ & $65$ & $66$ & $66$ & $69$\\
				\hline
				$71$ & $72$ & $72$ & $73$ & $73$ & $80$ & $81$ & $83$\\
				\hline
				$84$ & $85$ & $88$ & $89$ & $93$ & $94$ & $98$ & $265$\\
				\hline
			\end{tabular}
		\end{center}
		\begin{itemchoice} 
			\itemch Mỗi tháng cửa hàng bán trung bình $83{,}75$ bao.
			\itemch Số trung vị là $73$.
			\itemch Sai khác giữa số trung bình và số trung vị là $10{,}75$.
			\itemch Ta có số trung vị $Q_2=73$ .\\
			Số trung vị của nửa bên trái $Q_2$ là $Q_1=66$ .\\
			Khoảng cách từ $Q_1$ đến $Q_2$ là $7$.
		\end{itemchoice}
	}
\end{ex}

%%%========Câu 4
\begin{ex}%[0D6V3-4]
	Cho biết tổng diện tích rừng (đơn vị: triệu ha) từ năm $2008$ đến năm $2019$ ở nước ta lần lượt là $13{,}1$; $13{,}2$; $13{,}4$; $13{,}5$; $13{,}9$; $14{,}0$; $13{,}8$; $14{,}1$; $14{,}4$; $14{,}4$; $14{,}5$; $14{,}6$. Khi đó
	% \begin{center}
	% 	\begin{tikzpicture}
	% 		\matrix[matrix of nodes,nodes in empty cells,
	% 		row sep=-\pgflinewidth,column sep=-\pgflinewidth,
	% 		nodes={minimum height=8mm,minimum width=12.5mm,draw=cyan,anchor=center},
	% 		column 1/.style={nodes={minimum width=24mm,font=\bf,color=cyan}},
	% 		row 1/.style={nodes={fill=cyan!10}},
	% 		row 2/.style={nodes={minimum height=20mm}},
	% 		]{
	% 			Năm &2008&2009&2010&2011&2012&2013&2014&2015&2016&2017&2018&2019\\ 
	% 			\node[align=center]{Tổng diện\\tích rừng \\(triệu ha)}; &13{,}1&13{,}2&13{,}4&13{,}5&13{,}9&14,{,}0&13{,}8&14{,}1&14{,}4&14{,}4&14{,}5&14{,}6\\
	% 		};
	% 	\end{tikzpicture}
	% \end{center}
	\choiceTF
	{\True Diện tích rừng trung bình của nước ta từ năm $2008$ đến năm $2019$ khoảng $13{,}9$ triệu ha}
	{\True Từ năm $2008$ đến năm $2019$, diện tích rừng của năm có giá trị thấp nhất là $13{,}1$ triệu ha}
	{Từ năm $2008$ đến năm $2019$, diện tích rừng của năm có giá trị cao nhất là $14{,}0$ triệu ha}
	{\True So với năm $2008$, tỉ lệ tổng diện tích rừng của nước ta năm $2019$ tăng lên được $11{,}45\%$}
	\loigiai{	
		\begin{itemchoice}
			\itemch Diện tích rừng trung bình của nước ta từ năm $2008$ đến năm $2019$ là
			\begin{eqnarray*}
				S_{\text{tb}} &=& \dfrac{13{,}1 + 13{,}2 + 13{,}4 + 13{,}5 + 13{,}9 + 14{,}0 + 13{,}8 + 14{,}1 + 14{,}4 + 14{,}4 + 14{,}5 + 14{,}6}{12}\\
				&\approx&13{,}9~\text{(triệu ha)}.
			\end{eqnarray*}\vspace{-1cm}
			\itemch Năm $2008$ diện tích rừng có giá trị thấp nhất là $13{,}1$ (triệu ha).\\
			\itemch Năm $2019$ diện tích rừng có giá trị cao nhất là $14{,}6$ (triệu ha). 
			\itemch So với năm $2008$, tỉ lệ tổng diện tích rừng của nước ta năm $2019$ tăng lên
			$$\dfrac{14{,}6-13{,}1}{13{,}1}\cdot 100\% \approx 11{,}45\%.$$
		\end{itemchoice}
	}
\end{ex}

\Closesolutionfile{ans}
\indapan{3}{ans/ans-TF-C5-B3-D1}
\Opensolutionfile{ans}[ans/ans-KQ-C5-B3-D1]
\ind{PHẦN III.} \inden{Câu trả lời ngắn.}
\setcounter{ex}{0}
%%%%%%=========Câu 1
\begin{ex}%[0D6H3-3]
	\immini[thm]{
	Tìm trung vị của mẫu số liệu sau đây cho biết số bài tập Vật lí làm được trong một buổi học của học sinh lớp $10$
	}{
		\begin{tabular}{|c|c|c|c|c|c|}
			\hline
			Số bài & $1$ & $2$ & $3$ & $4$ & $5$ \\
			\hline
			Số học sinh & $4$ & $6$ & $9$ & $7$ & $4$ \\
			\hline
		\end{tabular}
		}
	\shortans[0]{3}
	\loigiai{
		Cỡ mẫu là $n=4+6+9+7+4=30$.\\
		Vì $n=30$ là số chẵn nên trung vị là trung bình của hai giá trị $x_{15}$  và  $x_{16}$.\\
		Ta có $x_{15}=x_{16}=3$.\\
		Vậy trung vị của mẫu số liệu là $Q_2=3$.
	}
\end{ex}

%%%%%%=========Câu 2
\begin{ex}%[0D6H3-4]
	\immini[thm]{Tìm tứ phân vị thứ nhất $Q_1$ của mẫu số liệu sau đây cho biết số giờ tự học Toán trong một tuần của học sinh lớp $10$}
	{\begin{tabular}{|c|c|c|c|c|c|}
			\hline
			Số giờ & $0$ & $1$ & $2$ & $3$ & $4$ \\
			\hline
			Số học sinh & $5$ & $7$ & $10$ & $6$ & $2$ \\
			\hline
		\end{tabular}}
	\shortans[0]{1}
	\loigiai{
		Cỡ mẫu là $n=5+7+10+6+2=30$.\\
		Vì $n=30$ là số chẵn nên $Q_1$ là trung vị của nửa dưới gồm $15$ số liệu đầu tiên.\\
		Vị trí của $Q_1$ là $x_8$.\\
		Ta thấy $x_8=1$, do đó $Q_1=1$.
	}
\end{ex}

%%%%%%=========Câu 3
\begin{ex}%[0D6V3-4]
	\immini[thm]{Bảng sau đây cho biết số lần học tiếng Anh trên internet trong một tuần của một học sinh lớp $10$.
	Hãy tìm các tứ phân vị cho mẫu số liệu này. Tính $T=Q_{1}+Q_{2}+Q_{3}$.}
	{\begin{tabular}{|l|l|l|l|l|l|l|}
			\hline
			Số lần & $0$ & $1$ & $2$ & $3$ & $4$ & $5$ \\
			\hline
			Số học sinh & $2$ & $4$ & $6$ & $12$ & $8$ & $3$ \\
			\hline
		\end{tabular}}
	\par
	\shortans[0]{9}
	\loigiai{
		Cỡ mẫu $n=2+4+6+12+8+3=35$.\\
		Vì cỡ mẫu là ${n=35}$, là số lẻ, nên giá trị của tứ phân vị thứ hai là ${Q_2=x_{18}=3}$.\\
		Tứ phân vị thứ nhất là trung vị của mẫu là ${x_1 ; x_2 ; \ldots ; x_{17}}$. Do đó ${Q_1=x_9=2}$.\\
		Tứ phân vị thứ ba là trung vị của mẫu ${x_{19} ; x_{20} ; \ldots ; x_{35}}$. Do đó ${Q_3=x_{27}=4}$.\\
		Vậy $T=Q_{1}+Q_{2}+Q_{3}=9$.
	}
\end{ex}



%%%%%%=========Câu 5
\begin{ex}%[0D6V3-5]
	Cho bảng phân bố tần số
	\begin{center}
		\newcolumntype{P}[1]{>{\centering\arraybackslash}p{#1}}
		\begin{tabular}{|c|P{0.6in}|P{0.6in}|P{0.6in}|P{0.6in}|P{0.6in}|}
			\hline
			Giá trị & $x_1$    & $x_2$    & $x_3$    & $x_4$    & $x_5$   \\
			\hline
			Tần số & 2    & $x+y$     & $2x-y$    & 5    & 6\\
			\hline
		\end{tabular}
	\end{center}
	với $x, y$ là các số tự nhiên. Có tất cả các cặp số $(x;y)$ để $x_5$ là mốt của bảng số liệu đã cho?
	\par
	\shortans{$14$}
	\loigiai{
		Điều kiện để $x_5$ là mốt của bảng số liệu đã cho là
		$\heva{ &0\le x+y\le 6& \quad (1)\\  &0\le 2x-y\le 6& \quad (2)\\}$\\
		Từ $(1)$ suy ra $ y\le 6-x $ và $ x\le 6-y $.\\
		Từ $(2)$ suy ra $ 2x-6\le y $ và $ \dfrac{y}{2}\le x $.\\
		Do đó $ \heva{&2x-6\le 6-x \\  &\dfrac{y}{2}\le 6-y \\ }\Leftrightarrow \heva{ &0\le x\le 4 \\  &0\le y\le 4 \\ }$\\
		Từ đó tìm được $ 14 $ cặp số   thỏa mãn là 
		$ \left(0;0 \right)$, $\left(1;0 \right)$, $\left(1;1 \right)$, $\left(1;2 \right)$, $\left(2;0 \right)$, $\left(2;1 \right)$, $\left(2;2 \right)$, 
		$\left(3;0 \right)$, $\left(3;1 \right)$, $\left(3;2 \right)$, $\left(3;3 \right)$, $\left(4;2 \right)$, $\left(2;3 \right)$, $\left(2;4 \right)$.
	}
\end{ex}


\Closesolutionfile{ans}
\indapan{6}{ans/ans-KQ-C5-B3-D1}
%%%%%%%%%%%%%%%%%%
\ind{PHẦN IV.} \inden{Tự luận.}
\setcounter{ex}{0}
%%%%%%=========Câu 1
\begin{ex}%[0D6H3-4]
	Tìm số trung bình, trung vị, mốt và tứ phân vị của mỗi mẫu số liệu sau đây:
	\begin{listEX}[1]
		\item Số điểm mà năm vận động viên bóng rổ ghi được trong một trận đấu
		\[9\quad 8\quad 15\quad 8\quad20\]
		\item Giá của một số loại giày (đơn vị nghìn đồng)
		\[350\quad 300\quad 650\quad 300\quad 450\quad 500\quad 300\quad 250\]
		\item Số kênh được chiếu của một số hãng truyền hình cáp
		\[36\quad 38\quad 33\quad 34\quad 32\quad 30\quad 34\quad 35\]
	\end{listEX}
	\loigiai{	
		\begin{listEX}[1]
			\item Sắp xếp lại mẫu số liệu $8\quad 8\quad 9\quad 15\quad20$.\\
			Số trung bình là $\overline{x}=\dfrac{8+8+9+15+20}{5}=12$.\\
			Trung vị là $Q=9$.\\
			Mốt là $M_o=8$.\\
			Tứ phân vị là $Q_1=\dfrac{8+8}{2}=8$, $Q_2=9$, $Q_1=\dfrac{15+20}{2}=17{,}5$.
			\item Sắp xếp lại mẫu số liệu $250\quad 300\quad 300\quad 300\quad 350\quad 450\quad 500\quad 650$.\\
			Số trung bình là $\overline{x}=\dfrac{250+300+300+300+350+450+500+650}{8}=387{,}5$.\\
			Trung vị là $Q=\dfrac{300+350}{2}=325$.\\
			Mốt là $M_o=300$.\\
			Tứ phân vị là $Q_1=300$, $Q_2=325$, $Q_1=\dfrac{450+500}{2}=475$.
			\item Sắp xếp lại mẫu số liệu $30\quad 32\quad 33\quad 34\quad 34\quad 35\quad 36\quad 38$.\\
			Số trung bình là $\overline{x}=\dfrac{30+32+33+34+34+35+36+38}{8}=34$.\\
			Trung vị là $Q=34$.\\
			Mốt là $M_o=34$.\\
			Tứ phân vị là $Q_1=\dfrac{32+33}{2}=32{,}5$, $Q_2=34$, $Q_1=\dfrac{35+36}{2}=35{,}5$.
		\end{listEX}
	}
\end{ex}

%%%%%%=========Câu 2
\begin{ex}%[0D6H3-4]
	Chiều cao (đơn vị: xăng-ti-mét) của các bạn tổ $I$ ở lớp $10A$ lần lượt là
	\[165\quad 155\quad 171\quad 167\quad 159\quad 175\quad 165\quad 160\quad 158\]
	Đối với mẫu số liệu trên, hãy tìm
	\begin{listEX}[4]
		\item Số trung bình cộng.
		\item Trung vị.
		\item Mốt.
		\item Tứ phân vị.
	\end{listEX}
	\loigiai{
		Mẫu số liệu trên được sắp xếp theo thứ tự tăng dần như sau
		\[155\quad 158\quad 159\quad 160\quad 165\quad 165\quad 167\quad 171\quad 175\]
		\begin{listEX}[1]
			\item Số trung bình cộng là
			$$\overline{x}=\dfrac{155+158+159+160+165+165+167+171+175}{9} \approx 163{,}9 \text{ (cm)}.$$
			\item Ta có $N=9$ là số lẻ. Số liệu thứ $\dfrac{9+1}{2}=5$. Vậy số trung vị là $M_e=165$.
			\item Ta thấy giá trị $165$ có tần số $2$ lớn nhất, do đó mốt của mẫu số liệu trên là $M_o=165$.
			\item Trung vị của dãy $155\quad 158\quad 159\quad 160$ là $Q_1=\dfrac{158+159}{2}=158{,}5$.\\
			Trung vị của dãy $165\quad 167\quad 171\quad 175$ là $Q_3=\dfrac{167+171}{2}=169$.\\
			Vậy $Q_1=158{,}5$ (cm), $Q_2=165$ (cm), $Q_3=169$.
		\end{listEX}
	}
\end{ex}



%%%%%%=========Câu 3
\begin{ex}%[0D6V3-5]
	Điểm bài kiểm tra một tiết môn toán của $40$ học sinh lớp $11A1$ được thống kê bằng bảng số liệu dưới đây
	\begin{center} 
			% {\renewcommand\arraystretch{1.5} %chỉnh chiều rộng dòng
				\begin{tabular}{|c|c|c|c|c|c|c|c|c|c|} 
					\hline 
					Điểm &  $3$ &  $4$ &  $5$ &  $6$ &  $7$ &  $8$ &  $9$ &  $10$ & Cộng \\
					\hline
					Số học sinh &  $2$ &  $3$ &  $3n-8$ &  $2n+4$ &  $3$ &  $2$ &  $4$ &  $5$ & $40$ \\
					\hline  
				\end{tabular}
		\end{center}
	Trong đó $n\in \mathbb{N},n\ge 4$. Tính mốt của bảng số liệu thống kê đã cho. 
	\loigiai{
		Vì tổng các số liệu thống kê bằng $40$ nên ta có $5n+15=40\Leftrightarrow n=5$.\\ 
		Với $n=5$ ta có bảng phân bố tần số
		\begin{center} 
			% {\renewcommand\arraystretch{1.5} %chỉnh chiều rộng dòng
				\begin{tabular}{|c|c|c|c|c|c|c|c|c|c|} 
					\hline 
					Điểm &  $3$ &  $4$ &  $5$ &  $6$ &  $7$ &  $8$ &  $9$ &  $10$ & Cộng \\
					\hline
					Số học sinh &  $2$ &  $3$ &  $7$ &  $14$ &  $3$ &  $2$ &  $4$ &  $5$ & $40$ \\
					\hline  
				\end{tabular} 
		\end{center}
		Vậy mốt của bảng số liệu là $M_o=6$.
	}
\end{ex}

%%%%%%=========Câu 4
\begin{ex}%[0D6V3-5]
	Cho bảng phân bố tần số sau
	\begin{center}
		\begin{tabular}{|c|c|c|c|c|c|}
			\hline
			Giá trị & $x_1$    & $x_2$    & $x_3$    & $x_4$    & $x_5$ \\
			\hline
			Tần số & 12    & 5     & $n^2$    & 16    & $6n-5$ \\
			\hline
		\end{tabular}
	\end{center}
	Tìm tất cả các số tự nhiên $n$ để $M_o=x_3$ là mốt duy nhất của bảng phân bố tần số đã cho.
	\loigiai{
		Từ giả thiết $M_o=x_3$ là mốt duy nhất của bảng số liệu thống kê đã cho nên ta có
		\[\begin{cases} n^2>16\\ n^2>6n-5\end{cases} 
		\Leftrightarrow \begin{cases} 
			\left[\begin{aligned}&n<-4\\&n>4\end{aligned}\right.\\
			\left[\begin{aligned}&n<1\\&n>5\end{aligned}\right. 
		\end{cases} \Leftrightarrow \left[\begin{aligned}&n<-4\\&n>5.\end{aligned}\right.\]
		Vì $n$ là số tự nhiên nên các giá trị $n$ thỏa mãn là $n>5$.
	}
\end{ex}


%%%%%%=========Câu 5
\begin{ex}%[0D6V3-5]
	Cho bảng phân bố tần số sau
	\begin{center}
		\begin{tabular}{|c|c|c|c|c|c|c|}
			\hline
			Giá trị & $x_1$    & $x_2$    & $x_3$    & $x_4$    & $x_5$   & $x_6$\\
			\hline
			Tần số & 5    & $n^2+3$     & 3    & $7n-9$    & $n+1$& 7\\
			\hline
		\end{tabular}
	\end{center}
	Gọi $S$ là tập hợp tất cả các số $n$ nguyên dương sao cho $M_o=x_2$ và $M_o=x_4$ là hai mốt của bảng phân bố tần số đã cho. Tính số phần tử của tập hợp $S$.
	\loigiai{
		Từ giả thiết $x_2$ và $x_4$ là các mốt của bảng số liệu thống kê đã cho, ta có
		$$\heva{& n^2+3=7n-9 \\ &7n-9>n+1 \\ &7n-9>7}
		\Leftrightarrow \heva{ &n^2-7n+12=0 \\ &n>\dfrac{5}{3}\\ &n>\dfrac{16}{7} }
		\Leftrightarrow \hoac{& n=3 \\ & n=4.}$$
		Vì $n$ là số nguyên dương nên $n=3$ và $n=4$ thỏa mãn. Vậy tập hợp $S$ có $2$ phần tử.
	}
\end{ex}

%%%%%%=========Câu 6
\begin{ex}%[0D6C3-2]
	Học sinh tỉnh $A$ (gồm lớp $11$ và lớp $12$) tham dự kì thi học sinh giỏi Toán của Tỉnh (thang điểm $20$) và điểm trung bình của họ là $10$. Biết rằng số học sinh lớp $11$ nhiều hơn số học sinh lớp $12$ là $50\%$ và điểm trung bình của khối $12$ cao hơn điểm trung bình của khối $11$ là $50 \%$. Điểm trung bình của khối $12$ là
	\loigiai{\begin{itemize}
			\item Gọi số học sinh lớp $12$ là $n$. Theo bài ra, số học sinh lớp $11$ sẽ là $1{,}5n$. Gọi điểm trung bình của học sinh lớp $11$ là $a$. Theo bài ra, điểm trung bình của học sinh lớp $12$ là $1{,}5 a$.
			\item Tổng số điểm của học sinh lớp $11$ là $S=a \cdot 1{,}5 n=1{,}5an$.
			\item Tổng số điểm của học sinh lớp $12$ là $T=(1{,}5a)n=1{,}5an$.
			Vậy tổng số điểm của học sinh lớp $11$ và $12$ là $1{,}5an+1{,}5 a n=3an$.
			\item Mặt, ta có tổng số học sinh lớp $11$ và $12$ là $n+1{,}5 n=2{,}5n$ và điểm trung bình của lớp $11$ và $12$ là $10$ . Do đó, tổng số điểm của học sinh lớp $11$ và $12$ là $10 \cdot(2{,}5n)=25n$.
			\item Từ đó ta có $3an=25n$ hay $a=\dfrac{25}{3}$,\\
			Vậy điểm trung bình của học sinh lớp 12 là $1{,}5a=1{,}5 \cdot \dfrac{25}{3}=12{,}5$.
		\end{itemize}
	}
\end{ex}
